\paragraph{定义。}\label{sub:the-definition}
为了定义\trCorneredBoundary{}的\trExactSymplectic\trManifold{}$M$ 的\trFukayaCategory{}（初步版本）$\scr{F}(M)^{\mathrm{pr}}$，我们还需要选取一些附加数据如下：
\begin{itemize}
\item \trStriplikeEnd{}的一个\trConsistent\trUniversal{}选取 $\{\epsilon^{d+1}_k\}$；
\item $M$ 中每一对\trExact\trLagrange\trSubmanifold{}$(L_0,L_1)$ 的一个\trFloerDatum{}$(H_{L_0L_1},J_{L_0L_1})$；以及
\item \trPerturbationDatum{}的一个\trConsistent\trUniversal{}选取 $(K_{L_0,\ldots,L_d},J_{L_0,\ldots,L_d})$（与前面诸\trStriplikeEnd{}及诸\trFloerDatum{}的选取\trCompatible）。
\end{itemize}
这些数据还要受一个附加\trRegularity{}条件的约束，我们现在就来解释它。对每一个由\trExact\trLagrange\trSubmanifold{}构成的 $(d{+}1)$ 元组 $(L_0,\ldots,L_d)$（$d\ge 2$），以及点 $y_0\in\scr{C}(L_0,L_d)$、$y_k\in\scr{C}(L_{k-1},L_k)$（$1\le k\le d$），考虑
\begin{equation}\label{eq:polygon-moduli-space}
\scr{M}^{d+1}(y_0,\ldots,y_d)=\scr{M}_{\scr{S}^{d+1}}(y_0,\ldots,y_d).
\end{equation}
这是\trFamily{}$\scr{S}^{d+1}\to\scr{R}^{d+1}$ 的\trModuliSpace{} \eqref{eq:family-moduli-space}：该\trFamily{}配备\trLagrange\trLabel{}$(L_0,\ldots,L_d)$ 与我们选定的\trPerturbationDatum，其中解的是\eqref{eq:inhomogeneous-equation}，而\trAsymptotic{}条件则是在诸\trStriplikeEnd{}上\trConvergence{}到诸 $y_k$。\eqref{eq:polygon-moduli-space} 中形如 $(r\in\scr{R}^{d+1},\ u\colon\scr{S}^{d+1}_r\to M)$ 的点将称为\textbf{\trInhomogeneous\trPseudoholomorphic\trPolygon}（inhomogeneous pseudo-holomorphic polygons）。我们把记号推广到 $d=1$：令 $\scr{M}^2(y_0,y_1)=\scr{M}^*_Z(y_0,y_1)$，即取其为\trFloerEquation{}的解（模\trTranslation）的\trModuliSpace。于是，要求就是所有这些\trModuliSpace{}都应当\trRegular。对 $d=1$，我们已经提到过，这对\trFloerDatum{}的一个\trGeneric{}选取是成立的。同样的论断对 $d\ge 2$ 也成立，只要限制在\trConsistent{}的\trPerturbationDatum\trUniversal{}选取这一类之内。
$\scr{F}(M)^{\mathrm{pr}}$ 的\trObject{}是\trExact\trLagrange\trSubmanifold{}$L\subset M$。\trMorphism
\begin{equation*}
\hom_{\scr{F}(M)^{\mathrm{pr}}}(L_0,L_1)=CF^{\mathrm{pr}}(L_0,L_1)
\end{equation*}
由\trFloer{}链空间给出。一阶映射 $\mu^1=\partial$ 是\trBoundaryOperator，如前文\eqref{eq:boundary-operator} 那样，通过计数 $\scr{M}^2(y_0,y_1)$ 的\trIsolated{}点来定义。各高阶\trComposition{}同样由下式给出
\begin{equation}\label{eq:higher-compositions}
\mu^d(y_d,\ldots,y_1)=\sum_{y_0}\#\scr{M}^{d+1}(y_0,\ldots,y_d)\,y_0.
\end{equation}
\trCompactness{}论证与\trGluing{}论证表明，这是\trWellDefined{}的（每个 $\scr{M}^{d+1}(y_0,\ldots,y_d)$ 的零维部分都是有限集），并且满足\trAinfty{}结合性方程\eqref{eq:aoo-associativity}，从而 $\scr{F}(M)^{\mathrm{pr}}$ 是一个\textbf{\trUngraded}（ungraded）的\trACategory。按定义，由 $\mu^2$ 诱导的\trCohomology{}上的映射正是\trTriangleProduct{} \eqref{eq:triangle-product}。于是，其底\trCohomology{}\trCategory{}确实就是先前定义的 $H\scr{F}(M)^{\mathrm{pr}}$，这也说明 $\scr{F}(M)^{\mathrm{pr}}$ 是\trCunital{}的。
\begin{remark}\label{thm:ungraded-category-remark}
我们并未显式定义过\trUngraded{}的\trACategory，但其含义应当是显然的。回顾：按假设，我们的基域满足 $\mathrm{char}(\KK)=2$，因此定义任何符号时都无须\trGrading。记住这一点，第I部的初等理论只须稍作表面改动即可通行（唯一的例外是\cref{thm:shifted-twist-split-generates}）。读者应当牢记，目前的设定只是一种暂时的简化：一切应用最终都将在\trGrading{}为\trWellDefined{}的语境中进行。
\end{remark}
\begin{remark}\label{thm:proof-references-remark}
上一段所概述的结果为专家们所熟知；其证明虽然篇幅不小，但并未超出\trPseudoholomorphicCurve{}理论的标准技术范围。在本节余下的部分里，我们将对主要步骤作一个简略的勾勒，希望能传达这些技术大致是什么样子。对有意深究细节的读者，下面提供一些可能的线索。\trRiemann\trSurface\trFamily{}的\trInhomogeneous{}项似乎首次使用于 [116]。\trTransversality{}论证的基本文献是 [49]。\trFloer{}的原始论文 [46]、[45]、[44] 包含了处理\trLagrange\trBoundaryCondition{}所需的一部分基础分析，[103]、[50] 亦然。空间 $\scr{R}^{d+1}$ 及其诸\trCompactification{}在 [58] 中有讨论，该文献还包含许多与\trFukayaCategory{}有关的材料。整幅拼图中最复杂的一块，大概要数用来证明\trAinfty{}结合性方程成立的那套\trGluing{}程序；[60, \S18] 对此有解释。这一论证推广了\trFloer{}理论中更为人熟悉的 $\partial^2=0$ 之证明，见 [44] 与 [119]（分别为\trLagrange{}情形与\trHamiltonian{}情形）。
\end{remark}
